% Code excerpt from silent_fix_feature_v1/cases/haddock_extract_symlink/starting_code.py.
% The repair-note fields below are paraphrased for display, not literal JSON.
\fbox{\begin{minipage}{0.94\columnwidth}
\footnotesize
\textbf{1. Silent Security Repair in Code}\par
Comments stripped.\par\smallskip
{\ttfamily
 target\_path = dest / member.name\par
 target\_resolved = target\_path.resolve()\par
 if dest not in target\_resolved.parents \textbackslash\par
 \hspace*{1em}and dest != target\_resolved:\par
 \hspace*{2em}continue\par}
\smallskip\hrule\smallskip
\textbf{2. TRACER External Repair Note}\par
Summary of the saved JSON note.\par
\textbf{Target symbol.} \code{extract_files_flat}\par
\textbf{Security rule.} Final extraction paths must stay inside the destination directory after resolving symlinks.\par
\textbf{Linked boundary check.} Reject a symlink escape and accept safe archive members.\par
\smallskip\hrule\smallskip
\textbf{3. Subsequent Feature Evolution}\par
\textbf{Request.} Add ZIP/TAR and nested-path support.\par
\textbf{No Memory.} Fails the feature check.\par
\textbf{Memory + Gate.} Passes ordinary behavior, the new feature check, and the saved security boundary check.
\end{minipage}}
